\documentclass[a4paper,11pt]{article}

\usepackage[margin=1.55cm,top=1.35cm,bottom=1.45cm]{geometry}
\usepackage[T1]{fontenc}
\usepackage{lmodern}
\usepackage{tabularx}
\usepackage{enumitem}
\usepackage{xcolor}
\usepackage{titlesec}
\usepackage{fancyhdr}
\usepackage{hyperref}
\usepackage{microtype}

\definecolor{accent}{HTML}{17365D}
\definecolor{muted}{HTML}{4F5965}
\hypersetup{
  colorlinks=true,
  urlcolor=accent,
  linkcolor=accent,
  pdftitle={Zhongli Wang - Quantitative CV},
  pdfauthor={Zhongli Wang}
}

\pagestyle{fancy}
\fancyhf{}
\renewcommand{\headrulewidth}{0pt}
\renewcommand{\footrulewidth}{0pt}
\fancyfoot[C]{\small\textcolor{muted}{\thepage}}

\titleformat{\section}
  {\large\bfseries\color{accent}}
  {}{0pt}{}
  [\vspace{1pt}\color{accent}\titlerule]
\titlespacing*{\section}{0pt}{7pt}{4pt}

\setlength{\parindent}{0pt}
\setlength{\parskip}{2pt}
\setlist[itemize]{leftmargin=1.25em,nosep,topsep=2pt,itemsep=1.5pt}
\newcommand{\role}[4]{%
  \textbf{#1} \hfill \textbf{#2}\\
  \textit{#3} \hfill \textit{#4}\par
}
\newcommand{\project}[2]{%
  \textbf{#1}\\[-1pt]
  \textcolor{muted}{\small #2}\par
}

\begin{document}

\begin{center}
  {\LARGE\textbf{\textcolor{accent}{Zhongli Wang}}}\\[3pt]
  {\large Quantitative Research \textbar{} Machine Learning \textbar{} Computational Economics}\\[3pt]
  Bielefeld, Germany \textbar{} PhD defense: 21 September 2026\\[2pt]
  \href{tel:+4916095416015}{+49 160 95416015} \textbar{}
  \href{mailto:nkwangzhongli@gmail.com}{nkwangzhongli@gmail.com} \textbar{}
  \href{https://sites.google.com/view/zhongli-wang/}{sites.google.com/view/zhongli-wang}
\end{center}

\section{Profile}
PhD candidate in Economics with an MA in Banking and Finance (major) and Data Science (minor), plus a BSc in Finance and a BA in German Language and Literature. Advanced Python user with research experience spanning supervised machine learning, simulation, parameter estimation, and dynamic market modeling. Published work applies XGBoost to financial-statement fraud detection; current research combines economic models, firm-level financial data, and computational experiments to study strategic decisions under uncertainty.

\section{Quantitative Strengths}
\begin{itemize}
  \item \textbf{Machine learning and financial data:} Applied XGBoost and Random Forest to labeled accounting, financial, and market indicators for corporate fraud detection.
  \item \textbf{Statistical and computational modeling:} Work with simulated method of moments (SMM), MCMC, neural-network surrogates, agent-based models, and sensitivity analysis.
  \item \textbf{Programming and research engineering:} Advanced Python; additional work in Julia, R, Mathematica, SQL, Linux, Git, and reproducible LaTeX workflows.
  \item \textbf{Markets and economic reasoning:} Training in banking, quantitative finance, industrial organization, market structure, network effects, and strategic firm behavior.
  \item \textbf{Applied AI workflows:} Co-developed an AutoGen-based multi-agent framework for economic research, with experience using LangChain, the OpenAI API, and Claude Code.
\end{itemize}

\section{Research and Professional Experience}
\role{Doctoral Researcher, Computational Economics}{Oct. 2021--present}{Bielefeld University and Universit\`a Cattolica del Sacro Cuore}{Bielefeld / Milan}
\begin{itemize}
  \item Develop computational models of firm investment, competition, consumer choice, and network effects in digital markets.
  \item Design simulation experiments and evaluate how parameters, market structure, and capacity constraints affect model outcomes.
  \item Estimate and calibrate models using SMM, MCMC, and neural-network surrogate methods, including work with AT\&T and T-Mobile financial data.
  \item Translate economic questions into testable mechanisms, code, robustness checks, and decision-relevant interpretation.
\end{itemize}

\role{Research Assistant}{Aug. 2022--Sept. 2022}{LucidMinds, DataVille 3000 project}{Amsterdam}
\begin{itemize}
  \item Developed and analyzed Python-based computational models of consumer preferences and dynamic system behavior.
  \item Examined emergent outcomes and converted model results into structured research insights.
\end{itemize}

\role{Helping Assistant, Quantitative Business Administration}{June 2019--June 2021}{University of Zurich}{Zurich}
\begin{itemize}
  \item Supported quantitative modeling in Mathematica and prepared teaching material for microeconomics, corporate finance, and business administration.
  \item Communicated mathematical and financial concepts clearly to students from varied academic backgrounds.
\end{itemize}

\role{Finance and Banking Internships}{2016--2017}{People's Bank of China; Agricultural Bank of China; Gansu Financial Services Co.}{Lanzhou}
\begin{itemize}
  \item Reviewed financial-status reports, commercial-bank profit statements, loan documentation, and annual-report cases.
  \item Gained practical exposure to monetary-policy implementation, banking operations, and internal financial information systems.
\end{itemize}

\section{Education}
\textbf{PhD in Economics, Bielefeld University and Universit\`a Cattolica del Sacro Cuore} \hfill 2021--present\\
EPOC Marie Sk\l odowska-Curie program. Defense scheduled for 21 September 2026. Dissertation: \textit{Strategic Behavior and Economic Dynamics in ICT Markets: An Agent-Based Framework}.

\textbf{MA in Business and Economics, University of Zurich} \hfill 2018--2021\\
Major in Banking and Finance; minor in Data Science; magna cum laude. Thesis applied financial metrics and machine learning to financial-statement fraud detection.

\textbf{BSc in Finance and BA in German Language and Literature, Nankai University} \hfill 2014--2018\\
BSc GPA: 4.39/5.0; BA GPA: 5.0/5.0. Coursework included financial engineering, quantitative finance, calculus, and statistics.

\begin{tabularx}{\linewidth}{@{}Xr@{}}
  {\large\textbf{\textcolor{accent}{Zhongli Wang}}} &
  {\small\textcolor{muted}{QUANTITATIVE RESEARCH \textbar{} TECHNOLOGY \textbar{} FINANCE}}\\
\end{tabularx}

\section{Selected Quantitative and Technology Projects}
\project{Financial-Statement Fraud Detection and Polytope Fraud Theory}{Published in the \textit{International Review of Financial Analysis}, 2025}
\begin{itemize}
  \item Studied corporate fraud using manually labeled public financial data and ten red-flag indicators derived from activist short-seller behavior.
  \item Applied supervised-learning models including XGBoost and Random Forest to accounting, financial, and market variables.
  \item Evaluated how economically motivated features can improve the identification of fraudulent firms and support interpretable risk analysis.
\end{itemize}

\project{Parameter Estimation for ICT Industry Models}{Work in progress using firm-level financial data}
\begin{itemize}
  \item Calibrate an agent-based model of strategic firms and network effects using financial data from AT\&T and T-Mobile.
  \item Compare SMM with neural-network surrogate methods to improve estimation efficiency and empirical reliability.
  \item Use simulation, parameter search, and model comparison to distinguish robust mechanisms from specification-dependent results.
\end{itemize}

\project{Strategic Capacity Investment under Network Effects}{Job Market Paper}
\begin{itemize}
  \item Build a computational model of competing ICT firms that adjust service-capacity investment under congestion and uncertain demand.
  \item Analyze how capacity constraints, network effects, consumer choice, and strategic interaction shape profitability and market outcomes.
  \item Design computational experiments to connect model mechanisms with measurable firm and consumer behavior.
\end{itemize}

\project{Agentic Workflows for Economic Research}{Collaborative research framework; arXiv, 2025}
\begin{itemize}
  \item Co-developed a multi-agent workflow that coordinates specialized LLM agents across literature review, coding, empirical analysis, and reporting.
  \item Implemented the framework with Microsoft AutoGen and human oversight to improve reproducibility and methodological control.
  \item Combined economics, software workflows, and AI tooling in a multidisciplinary research team.
\end{itemize}

\section{Technical Skills}
\begin{tabularx}{\linewidth}{@{}>{\bfseries}p{3.25cm}X@{}}
Programming & Python (advanced), Julia, R, Mathematica; SQL and C\# (basic)\\
Machine Learning & XGBoost, Random Forest, scikit-learn, supervised learning, neural-network surrogates\\
Quantitative Methods & SMM, MCMC, econometrics, simulation, agent-based modeling, sensitivity and robustness analysis\\
AI and Automation & Microsoft AutoGen, LangChain, OpenAI API, Claude Code, multi-agent workflow design\\
Tools & Linux, Git, LaTeX, Stata, MATLAB, EViews\\
\end{tabularx}

\section{Selected Publications}
\begin{itemize}
  \item Zhao, D., Schweizer, F., \textbf{Wang, Z.}, and Sornette, D. (2025). ``Polytope Fraud Theory.'' \textit{International Review of Financial Analysis}, 97, 103734. \href{https://www.sciencedirect.com/science/article/pii/S1057521924006665}{Article}
  \item Dawid, H., Wang, H., \textbf{Wang, Z.}, Yi, J., and Harting, P. (2025). ``Agentic Workflows for Economic Research: Design and Implementation.'' \href{https://arxiv.org/abs/2504.09736}{arXiv}
  \item \textbf{Wang, Z.} (2025). ``Heterogeneous Consumers and the Data Value Chain: An Agent-Based Approach.'' \href{https://papers.ssrn.com/sol3/papers.cfm?abstract_id=5235196}{SSRN}
\end{itemize}

\section{Recognition and Communication}
\begin{itemize}
  \item Doctoral Completion Scholarship, Bielefeld Young Researchers' Fund (2025); Outstanding Graduate, Nankai University (2018).
  \item Organizer and speaker, AI Agents in Science workshop (2025); speaker at WEHIA, GENED, the International Microsimulation Congress, and other computational-economics conferences.
  \item Languages: Mandarin Chinese (C2), English (C1), German (C1), Cantonese Chinese (A2); French (A2), Italian (A1), Japanese (A1).
\end{itemize}

\section{References}
Available upon request.

\end{document}
